

The corresponding metrics are detailed in \autoref{tab:ipw_diag_summary}.

\begin{table}[H]
\centering
\caption{Development-Occurrence IPW Diagnostic Summary (Stabilized and Truncated Weights)}
\label{tab:ipw_diag_summary}
\begin{tabular}{lccc}
\toprule
\textbf{Specification} & \textbf{Weight Rule} & \textbf{ESS} & \textbf{ESS / N} \\
\midrule
Unweighted reference & $w_i \equiv 1$ & 7,074 & 100.0\% \\
Stabilized (no truncation) & $\Pr(D=1)/\hat e_i$ & 412 & 5.8\% \\
Primary truncation & $\Pr(D=1)/\hat e_i$, clipped to $[p_{1},p_{99}]$ & 1,268 & 17.9\% \\
Sensitivity band A & clipped to $[p_{2.5},p_{97.5}]$ & 1,611 & 22.8\% \\
Sensitivity band B & clipped to $[p_{5},p_{95}]$ & 1,985 & 28.1\% \\
\bottomrule
\multicolumn{4}{l}{\footnotesize Notes: $\hat e_i = \Pr(D=1\mid X_i)$ from the development-occurrence selection model. ESS computed as $ (\sum_i w_i)^2 / \sum_i w_i^2 $.} \\
\multicolumn{4}{l}{\footnotesize Diagnostic gate for acceptable weighting requires ESS/N $\ge 40\%$ plus balance criteria; all tested bands fail the ESS gate.} \\
\end{tabular}
\end{table}
